\documentclass[../main.tex]{subfiles}
\begin{document}
%===================TIÊU ĐỀ TẬP=============================
\begin{table}[h]
  \begin{adjustwidth}{-1cm}{}
    \begin{tabular}{>{\centering\arraybackslash}p{6cm}|>{\raggedright\arraybackslash}p{15cm}}
      \multirow{5}{6cm}
      % Cột bên trái: ÉPISODE và số tập
      {\\[-40pt]\flaregothic\fontsize{20pt}{20pt}\selectfont \centering
        \color{eptype}ÉPISODE\color{black}\\
        \burbank\fontsize{72pt}{72pt}\selectfont
        \color{epnum}-8\color{black}\\[6pt]
        \cafeteria\fontsize{20pt}{20pt}\selectfont\color{epnum}(P-3)\color{black}\\
      }
       &
      % Dòng thứ nhất: tên tập tiếng Nhật
      {\fotskip\fontsize{18pt}{18pt}\selectfont \ruby{神}{かみ}の\ruby{年末}{ねんまつ}
      }                         \\[6pt]
      % Dòng thứ hai: tên tập tiếng Anh
       & {\cafeteria\fontsize{18pt}{18pt}\selectfont God's End of the Year }                        \\[4pt]
      % Dòng thứ ba: tên tập tiếng Việt
       & {\vnmsans\fontsize{18pt}{18pt}\selectfont Cuối năm cùng với Thần}                          \\[8pt]
      % Dòng thứ tư: Nhân vật xuất hiện
       & {\caviar{\fontsize{14pt}{14pt}\selectfont Nhân vật: \emp{miko} \emp{kintoki} \emp{ghost}}} \\[8pt]
      % Dòng cuối cùng: Thời gian diễn ra của tập (thường sẽ là ngày phát hành)
       & {\continuum\fontsize{16pt}{16pt}\selectfont Ngày phát hành: 21/12/2018}
    \end{tabular}
  \end{adjustwidth}
\end{table}
%========================VĂN BẢN==============================
\color{black}
\txtbx[2]{
  {\color{vol} \uline{Tóm tắt:}} Những ngày cuối năm bận rộn đã đến, nhưng Ngôi đền Ảo lại rơi vào tình trạng báo động đỏ khi vị Thần linh cai quản nơi đây đột nhiên đình công và bỏ nhà ra đi. Lý do thì ai cũng biết: vì cô nàng Vu nữ Sakura Miko quá lười biếng nên đã đẩy hết mọi công việc lên đầu ông Thần tội nghiệp.
}

Bầu không khí bên trong căn phòng nghỉ quen thuộc của Miko hôm nay bỗng chốc trở nên hỗn loạn đến mức khó tin. Tiếng chuông báo động đỏ vang lên inh ỏi, chói tai, xé toạc đi sự yên tĩnh thanh bình thường ngày của Ngôi đền Ảo. Đứng giữa căn phòng đang nhấp nháy ánh đèn đỏ rực, cả Miko và Kintoki đều trố mắt sững sờ, chưa kịp tiêu hóa nổi tình huống thót tim đang diễn ra.

``C-Cái quái gì đang xảy ra vậy!?'' Cô nàng tóc hồng thảng thốt hét lên, hai tay ôm lấy đầu trong cơn hoảng loạn.

Đứng bên cạnh, Kintoki nhăn nhó mặt mày, bộ lông màu hồng dựng đứng cả lên. Chú mèo lầm bầm bằng một chất giọng run rẩy ngập tràn sự lo lắng: ``Miko\dots\ Lần này thì chúng ta nguy to thật rồi!''

\begin{figure}[H]
  \centering
  \begin{subfigure}[t]{0.45\textwidth}
    \centering
    \includegraphics[width=8cm]{../../../../Image/BookP/p1-3a1.png}
  \end{subfigure}
  \quad
  \begin{subfigure}[t]{0.45\textwidth}
    \centering
    \includegraphics[width=8cm]{../../../../Image/BookP/p1-3a2.png}
  \end{subfigure}
\end{figure}

Nguyên nhân dẫn đến sự hỗn loạn kinh thiên động địa này là gì? Mọi chuyện bắt nguồn từ một bức thư bí ẩn được ai đó nặc danh đặt ngay ngắn trên chiếc bàn uống trà giữa phòng. Miko run rẩy cầm lấy tờ giấy, đôi mắt đảo nhanh qua từng dòng chữ, và chỉ trong tích tắc, sắc mặt cô nàng đã chuyển từ hồng hào sang trắng bệch hệt như tờ giấy cô đang cầm.

\txtbx{
  \textit{``Ta thực sự không thể chịu đựng thêm được nữa rồi.}

  \textit{Kẻ làm chủ cái ngôi đền này dạo gần đây làm việc quá đỗi lơ là. Cô ta lười biếng, bỏ mặc mọi công việc tế lễ ở đền chùa, để ta phải một thân một mình gồng gánh cáng đáng hết mọi thứ. Kết quả là ta đã kiệt sức, không thể nghỉ ngơi và cũng chẳng còn chốn nào để về nữa.}

  \textit{Ta sẽ bỏ đi và không bao giờ quay lại cái ngôi đền bóc lột sức lao động này nữa đâu!}\\[6pt]
  \textit{Đừng tốn công vô ích đi tìm ta làm gì.''}

  \hfill -- Từ vị Thần Ảo vô danh đáng thương trong ngôi đền này.
}

\ct

Đọc xong bức thư đẫm nước mắt kia, Miko ngả lưng đánh phịch xuống ghế, đưa tay lên trán vỗ vỗ mấy cái rồi thở dài ngao ngán than vãn: ``Cuối cùng thì cũng đã đến ngày này rồi à\dots\ Năm nào cũng thế, cứ đến cuối năm là ông Thần ấy lại giở chứng than vãn khó nhọc thế này, ai mà chịu cho nổi chứ\dots'' Nhìn bộ dạng dửng dưng của cô nàng, dường như chuyện ``Thần linh bỏ nhà đi bụi'' đã trở thành một thủ tục thường niên mà cô đã quá quen thuộc.

Ngược lại với thái độ dửng dưng của cô chủ, Kintoki lúc này có vẻ không thể nào giữ nổi bình tĩnh. Chú mèo bay nhảy loạn xạ quanh phòng, vội vàng lên tiếng cảnh báo: ``Nếu cứ để tình trạng mất cân bằng thế này tiếp diễn, thì chẳng mấy chốc cả cái thế giới Ngôi đền Ảo này sẽ loạn cào cào cho mà xem\dots\ thậm chí có khi còn sụp đổ và biến mất hoàn toàn không chừng!''

Nghe lời hù dọa của thức thần, nàng Vu nữ chỉ đành tặc lưỡi thở dài thêm một cái nữa. Cô lững thững lê bước đi vòng quanh phòng, miệng lầm bầm: ``Tôi nghĩ chắc ông ta lại đang bày trò hờn dỗi trẻ con gì ở loanh quanh đây thôi mà\dots''

``Chắc là trốn ở trong này nhỉ?'' Miko dừng bước trước chiếc tủ quần áo bằng gỗ nằm lặng lẽ ở một góc phòng. Bằng trực giác sắc bén của một ``Vu nữ ưu tú'', cô tin chắc mười mươi rằng vị Thần Ảo kia đang cố thủ ở trong đó.

\img{p1-3b}

Cô đưa tay gõ lộc cộc vào cánh cửa tủ gỗ, cất giọng khá bực mình pha chút thiếu kiên nhẫn: ``Thần ơi, ông có đang trốn ở trong đó không hả?''

Cánh cửa tủ vẫn im lìm đóng chặt, không hề có lấy một tiếng động đáp lại. Thế nhưng, ngay giây tiếp theo, một khung chữ nhật mờ ảo mang phong cách giao diện của các trò chơi visual novel (tiểu thuyết tương tác) bất thình lình hiện ra lơ lửng ngay trước cánh tủ. Bên dưới khung là những dòng chữ chạy rẹt rẹt mà có lẽ trên đời này, ngoài ``Vu nữ Ưu tú'' Miko và thức thần Kintoki ra thì chẳng có một ai khác có khả năng đọc hiểu được:

\txtbx{``Không có ai ở đây cả nha.''}

Miko nhíu mày, quay sang nhún vai đắc ý với Kintoki: ``Thấy chưa, tôi biết ngay tắp lự là ổng trốn trong này mà, ne!''

Chú mèo hồng bực dọc gắt gỏng hướng về phía cái tủ: ``Cái ông Thần kia, ông đừng có ở đó mà tỏ vẻ hờn dỗi như trẻ con lên ba nữa đi! Mau bước ra ngoài này làm việc!''

Thế nhưng, Vị Thần nọ dường như quyết tâm ăn vạ đến cùng. Trên cái khung chữ lơ lửng kia lại tiếp tục hiện ra những dòng tâm sự não nề ngập mùi nhõng nhẽo:

\txtbx{``Không chịu đâu\dots\ Ta mệt mỏi lắm rồi, ta không muốn làm việc nữa\dots\ Nhưng mà\dots\ ta lại vẫn muốn được ở bên cạnh Miko cơ\dots\ Sao cái số làm Thần linh của ta nó lại khổ ải thế này\dots''}

Cách hành xử khó hiểu và đầy mâu thuẫn này của đấng tối cao khiến nàng Vu nữ phải đưa tay ôm trán thở dài đầy bất lực. Cô thầm tự nhủ, nếu không giải quyết dứt điểm cái cục nợ này ngay bây giờ, có lẽ cô sẽ chẳng bao giờ có cơ hội được yên thân mà lười biếng nữa.

``Thôi được rồi, được rồi! Nếu ông thực sự muốn được tôi dỗ dành xoa dịu đến thế, thì để tôi chiều ông một bữa vậy.''

\ct

Nói là làm, cô nàng tóc hồng nắm chặt cán cây gậy \textit{gohei} đuổi tà quen thuộc của mình, bước từng bước chậm rãi tiến sát tới gần chiếc tủ - nơi ``ông Thần'' vẫn đang ngoan cố tử thủ bên trong, nhất quyết không chịu lòi mặt ra ngoài.

\img{p1-3c}

Miko hắng giọng, cố gắng ép tông giọng của mình trầm xuống thật ngọt ngào và tràn đầy sức thuyết phục: ``Ông thần ơi à\dots\ Trong suốt một năm qua, ông đã vất vả làm lụng rất nhiều rồi. Ông chỉ cần cố gắng kiên nhẫn thêm một chút xíu, chút xíu nữa thôi\dots\ Cố lên, ông Thần giỏi giang nhất trần đời ơi!''

Kèm theo những lời dỗ ngọt êm tai ấy, một vầng hào quang chói lọi bỗng chốc tỏa ra xung quanh người cô nàng. Cái dáng vẻ lúc này của Miko trông hệt như một người mẹ hiền hậu đầy bao dung, đang đặt tay lên vai vỗ về đứa con nhỏ đang mệt mỏi và kiệt sức của mình để tiếp thêm động lực bước tiếp.

Đứng ở một góc phòng, Kintoki chỉ biết thu mình lại, lơ lửng im lặng đứng nhìn tấn bi hài kịch giữa ``cô chủ'' của nó và ông Thần Ảo đang diễn ra trước mắt. Chú mèo đảo mắt, khẽ hừ lạnh một tiếng rồi lẩm bẩm trong miệng: ``Năm nay chắc mẩm cô ả lại đang diễn trò mỹ nữ kế đây mà\dots''

\img{p1-3d}

Và đúng như dự đoán, chỉ được đúng ba giây đóng vai hiền thục, bản chất thật của Miko đã lộ diện. Ngay lập tức, cô nàng gằn giọng, lớn tiếng mắng mỏ như thể vừa hoàn thành xong nghĩa vụ dỗ dành cho có lệ: ``Đó! Năm nay là tôi đã cực kỳ cất công dỗ dành an ủi ông một cách đường đường chính chính rồi nha! Liệu hồn mà biết điều, đừng có ở đó mà than vãn ỉ ôi nữa nhe!''

Dù cho lời nói đằng sau có cục súc đến đâu, thì có vẻ như những lời ngọt ngào ``hiếm có khó tìm'' trước đó của Miko đã phát huy tác dụng. Hai cánh cửa tủ gỗ (mà thực chất tượng trưng cho hai bên má của vị Thần đang nấp bên trong) bỗng chốc ửng đỏ hỏn lên vì ngượng. Vị Thần uy quyền giờ đây có vẻ đang rơi vào trạng thái xấu hổ tột độ, ngượng ngùng đến mức không biết nên xử lý tình huống này ra sao.

Dường như chính nàng Vu nữ cũng cảm thấy nổi da gà vì màn diễn sâu vừa rồi của mình, cô ả gắt lên, vung vẩy cây gậy trừ tà trên tay đập đập vào cánh cửa tủ: ``\textbf{Xong rồi thì nhanh nhanh cái chân lên và chui ra ngoài này đi! Làm việc tiếp thôi!}''

\img{p1-3e}

Cuối cùng, không thể chống lại sức ép quá lớn từ cả lời ngọt ngào lẫn sự đe dọa vũ lực của cô nàng Vu nữ, ``ông Thần Ảo'' cũng đành ngậm ngùi giương cờ trắng đầu hàng. Cánh cửa tủ từ từ hé mở, vị thần lềnh phềnh bay ra ngoài, chính thức chấm dứt trò chơi trốn tìm trẻ con này để quay lại với núi công việc dang dở.

Vậy là lại một ngày đầy những sự kiện dở khóc dở cười nữa khép lại trong Ngôi đền Ảo. Năm cũ bận rộn dần trôi qua, nhường chỗ cho những điều mới mẻ và thú vị đang chờ đón trong năm mới ở phía trước.

\end{document}